\chapter{Явная коэффициентная карта Тёплица для класса \texorpdfstring{$1/7$}{1/7}}

Литературный аудит установил общий механизм. Теперь он применяется к
конкретным данным проекта без дальнейшего поиска аналогий. Цель главы ---
построить граничный оператор, вычислить его дефектный проектор и проверить
совместимость двух ориентаций с KO6.

\section{Коэффициентный проектор}

В сравнительной алгебре
\begin{equation}
 M_{35}\otimes M_3\simeq M_{105}
\end{equation}
уже получен проектор
\begin{equation}
 q_0=p_{15}\otimes P_0,
 \qquad
 \rank q_0=15,
 \qquad
 \tau_{105}(q_0)=\frac{15}{105}=\frac17.
 \label{eq:v5-toeplitz-q0}
\end{equation}
Он каноничен относительно двух ранее выведенных данных: наблюдаемого угла
$p_{15}$ и голономного спектрального проектора $P_0$.

\section{Оператор на пространстве Харди}

Пусть $H^2(S^1)$ --- пространство неотрицательных фурье-мод на
ориентированной граничной окружности, а
\begin{equation}
 S e_n=e_{n+1},
 \qquad n=0,1,2,\ldots
\end{equation}
--- односторонний сдвиг. Для хопфова clutching-символа
\begin{equation}
 u_H(z)=z,
 \qquad \operatorname{wind}(u_H)=1,
\end{equation}
оператор Тёплица равен $T_{u_H}=S$.

На коэффициентном модуле определим
\begin{equation}
 \mathcal H_+=H^2(S^1)\otimes q_0\mathbb C^{105},
 \qquad
 \mathbb T_+=S\otimes q_0.
 \label{eq:v5-toeplitz-coefficient-operator}
\end{equation}
Тогда
\begin{equation}
 \mathbb T_+^*\mathbb T_+=I,
 \qquad
 I-\mathbb T_+\mathbb T_+^*
 =|e_0\rangle\langle e_0|\otimes q_0.
 \label{eq:v5-toeplitz-defect-projection}
\end{equation}
Правая часть является положительным компактным проектором ранга
пятнадцать. Следовательно,
\begin{equation}
 \operatorname{ind}\mathbb T_+=-15
\end{equation}
при стандартной конвенции $\dim\ker T-\dim\ker T^*$.

Композиция обычного следа на компактном ранг-один проекторе с
нормированным коэффициентным следом даёт
\begin{equation}
 \Tr_{\mathcal K}(|e_0\rangle\langle e_0|)
 \tau_{105}(q_0)=1\cdot\frac17=\frac17.
 \label{eq:v5-toeplitz-defect-weight}
\end{equation}

\begin{theorem}[Явный положительный представитель класса $1/7$]
Коэффициентный оператор Тёплица имеет компактный дефектный проектор
$|e_0\rangle\langle e_0|\otimes q_0$ ранга пятнадцать. Его нормированный
коэффициентный вес равен $1/7$, а $K_0$-класс совпадает по стабильному
рангу с $([p_{20}]-[p_{15}])\otimes[I_3]$.
\end{theorem}

Тем самым положительный представитель найден не как искусственный
проектор ранга пять внутри $H_{20}$, а в компактном идеале естественного
граничного расширения.

\section{Обращение ориентации и KO6}

Для сопряжённой линии $L^*$ символ равен
\begin{equation}
 u_H^*(z)=z^{-1}.
\end{equation}
Соответствующий оператор есть обратный сдвиг:
\begin{equation}
 \mathbb T_-=S^*\otimes q_0,
 \qquad
 \operatorname{ind}\mathbb T_-=+15.
\end{equation}
Вещественная структура обращает хопфову ориентацию и меняет два
оператора местами:
\begin{equation}
 J\mathbb T_+J^{-1}=\mathbb T_-.
 \label{eq:v5-toeplitz-J-pair}
\end{equation}
Поэтому полный сопряжённый пакет имеет нулевой обычный суммарный индекс,
тогда как каждая ориентированная половина хранит абсолютный индекс
пятнадцать с противоположным знаком. Это соответствует прежнему правилу
$E/E^*\mapsto L/L^*$ и не требует ручного выбора знака.

Однако данная проверка пока выполнена на уровне парного комплексного
индекса. Полный вещественный неограниченный $KKO$-цикл с правильными
KO6-знаками ещё не построен.

\section{Почему конечная матрица не воспроизводит результат}

Если заменить $S$ конечной $N\times N$ матрицей сдвига, она имеет
одномерные ядро и коядро, поэтому её индекс равен нулю. Нетривиальный
индекс возникает только в бесконечномерном модуле Харди:
\begin{equation}
 \operatorname{ind}S=-1,
 \qquad
 \operatorname{ind}S_N=0.
\end{equation}
Это не численная трудность, а принципиальное различие между конечной
матрицей и фредгольмовым оператором.

\begin{nogocriterion}[Запрет конечномерной имитации индекса]
Нельзя получить класс Тёплица заменой пространства Харди большим, но
конечным числом мод. Любая такая усечённая матрица имеет нулевой индекс и
не представляет граничную карту.
\end{nogocriterion}

\section{Является ли пространство Харди новым физическим сектором}

Пространство Харди здесь используется как аналитический представитель
класса расширения, а не как дополнительное множество внутренних
фермионов. $K$-класс не зависит от конкретного выбора базиса и устойчив
при компактных деформациях поляризации.

Тем не менее конечный родитель $M_{35}$ сам по себе не является алгеброй
Тёплица. Для физического включения требуется показать, что граничная
окружность дефекта, её ориентация и спектральное разложение действительно
являются граничными данными одного родительского спектрального тройного
объекта. Пока это не доказано, модуль Харди остаётся допустимым
индексным инструментом, но не физическим носителем.

Отдельно важно, что хопфова окружность в этой конструкции является
clutching-границей двух локальных карт сферы вокруг дефекта. Её нельзя
без доказательства отождествлять с компактным множителем $S^1$ пространства
$\mathbb{RP}^3\times S^1$, где ранее использовались иные спин-граничные
условия.

\section{Статус моста}

На комплексном $K$-теоретическом уровне мост теперь замкнут:
\begin{equation}
 [u_H]\boxtimes[q_0]
 \xmapsto{\partial_{\mathcal T}}
 \pm[|e_0\rangle\langle e_0|\otimes q_0]
 =\pm([p_{20}]-[p_{15}])\otimes[I_3].
 \label{eq:v5-toeplitz-closed-class-bridge}
\end{equation}
Знак фиксируется ориентацией $L/L^*$, положительный представитель явен,
а произвольность $\operatorname{Gr}(5,20)$ и $U(15)$ больше не входит в
определение класса.

Но это ещё не замыкание физической теории. Необходимо встроить расширение
в реальный KO6-родитель, получить неограниченный оператор и доказать, что
его компактный дефект действительно читается как локализованный
физический канал, а не только как индексная бухгалтерия.

\begin{protocol}
\begin{enumerate}
 \item коэффициентный проектор $q_0$: \textbf{каноничен};
 \item оператор Тёплица: \textbf{построен};
 \item компактный положительный проектор ранга пятнадцать:
 \textbf{получен};
 \item его вес: \textbf{$1/7$};
 \item совпадение со стабильным угловым классом: \textbf{получено};
 \item необходимость ручного $r_5\le p_{20}$: \textbf{устранена на
 уровне $K_0$};
 \item сопряжённые индексы KO6-пары: \textbf{$-15/+15$};
 \item суммарный обычный индекс полного пакета: \textbf{ноль};
 \item полный вещественный $KKO$-цикл: \textbf{не построен};
 \item пространство Харди как физический сектор: \textbf{не объявляется};
 \item комплексный классовый мост: \textbf{замкнут};
 \item физическое замыкание: \textbf{не достигнуто};
 \item следующий гейт: \textbf{вещественный KO6-подъём расширения}.
\end{enumerate}
\end{protocol}

Машинный сертификат сохранён в
\path{s2t_v5_one_seventh_toeplitz_boundary_map_gate_results.json}.